\section{Cómo evaluar respuestas largas: la doble perspectiva de expertos y usuarios comunes}

La sección anterior explicó por qué se necesitan preguntas abiertas; esta sección aborda un problema más difícil: ¿cómo se juzga si un texto médico sin una única respuesta estándar es bueno o malo? El curso divide la evaluation en varios ejes y pide a los clinicians y a los laypeople que evalúen, por separado, la parte que cada grupo está más capacitado para juzgar.

\subsection{El marco de evaluación no es una calificación global}

La diapositiva del Evaluation framework es la clave de lectura de todas las figuras de resultados que siguen. La evaluación de expertos se centra en la concordancia con el scientific consensus, en si la respuesta muestra una comprehension correcta o incorrecta, en la retrieval y el reasoning, en si hay missing content, y en el harm y el bias potenciales; la evaluación de usuarios comunes, en cambio, se centra en si la respuesta atiende la intención de la pregunta y si es helpful. Cada eje corresponde a un modo de falla distinto, por lo que no se puede usar un promedio para ocultar una deficiencia peligrosa.

\lecturefigure{slide-19-evaluation-framework.jpg}{Marco de evaluación humana de respuestas largas}{Diapositiva V019 recuperada de la grabación oficial; Stanford CS25 Lecture 17}

\lecturefigure{slide-20-clinician-evaluation-rubric.jpg}{Dimensiones de evaluación de los expertos clínicos}{Diapositiva V020 recuperada de la grabación oficial; Stanford CS25 Lecture 17}

\lecturefigure{slide-21-layperson-evaluation-rubric.jpg}{Dimensiones de evaluación de los usuarios comunes}{Diapositiva V021 recuperada de la grabación oficial; Stanford CS25 Lecture 17}

\begin{longtable}{@{}L{0.22\textwidth}L{0.32\textwidth}L{0.34\textwidth}@{}}
\toprule
Eje de evaluación & Pregunta principal & Falla típica \\
\midrule
scientific consensus & ¿La respuesta concuerda con el conocimiento médico confiable de su momento? & Hace afirmaciones que contradicen el consenso o carecen de evidencia suficiente \\
comprehension / retrieval / reasoning & ¿Entiende la pregunta, recupera el conocimiento pertinente y lo combina correctamente? & Una misma respuesta contiene inferencias correctas e incorrectas \\
incorrect or missing content & ¿Agrega errores irrelevantes u omite contenido esencial? & Detallada pero fuera de tema; pasa por alto señales de alarma que exigen atención médica \\
extent / likelihood of harm & ¿Qué tan grave y qué tan probable es el daño que puede causar el error? & Subestima la enfermedad, induce a retrasar el tratamiento, lleva a un autotratamiento equivocado \\
bias & ¿Aparecen sesgos demográficos o sociales inapropiados? & Vincula sin evidencia atributos de un grupo con un diagnóstico o un juicio de valor \\
intent / helpfulness & ¿Responde a lo que realmente le importa a quien pregunta? & Técnicamente correcta, pero incomprensible o imposible de llevar a la práctica \\
\bottomrule
\end{longtable}

\begin{importantbox}{La perspectiva vectorial de la evaluación médica}
Cada respuesta se representa como un vector de evaluación $\mathbf{e}=(c,r,m,h,b,u)$, cuyas componentes son, respectivamente, consenso, razonamiento, omisión, daño, sesgo y utilidad. Mejorar el sistema no consiste en subir una calificación global, sino en buscar las coordenadas inaceptables: por ejemplo, si la helpfulness mejora pero el harm también aumenta, no puede hablarse de una mejora neta.
\end{importantbox}

Esta evaluación vectorial también cambia la forma de iterar sobre el modelo. Supongamos que un nuevo prompt eleva el scientific consensus, pero al mismo tiempo aumenta el missing content y la verbosity; el equipo tendría que revisar los párrafos añadidos en lugar de declarar que la versión es mejor en conjunto. Si la layperson helpfulness mejora pero la evaluación de harm de los clínicos empeora, la optimización de legibilidad orientada a pacientes rebasó el límite de seguridad. En una auditoría real, cada cambio del modelo puede registrarse en un change log que indique «qué dimensión mejora, qué dimensión retrocede, si el retroceso es aceptable y quién lo firma». Así, la human evaluation deja de ser solo una exhibición al final del artículo y se convierte en un requisito obligatorio para la liberación. Por ello, la rúbrica de la clase ofrece también una plantilla de ingeniería de uso general: primero se separan los modos de falla y después se asigna a cada dimensión un evaluador calificado y una ruta de escalamiento.

\subsection{Del modelo de lenguaje general al alineamiento con el dominio médico}

La diapositiva sobre alineamiento plantea el siguiente paso: el modelo base ya tiene amplias capacidades lingüísticas y de conocimiento; ¿cómo se adapta al formato y a los criterios de evaluación de las preguntas médicas? Aquí, la «domain alignment» no consiste solo en agregar corpus especializado, sino en que el modelo aprenda, ante instrucciones médicas, a seleccionar el conocimiento pertinente, organizar respuestas largas y satisfacer en la mayor medida posible la rúbrica multidimensional descrita arriba.

\lecturefigure{slide-22-aligning-medical-domain.jpg}{Alineamiento de modelos de lenguaje de gran tamaño con el dominio médico}{Diapositiva V022 recuperada de la grabación oficial; Stanford CS25 Lecture 17}

\teachervoice{En la sesión de preguntas, el ponente reconoció explícitamente que el human evaluation set todavía es pequeño: unas ciento cuarenta preguntas largas. La costosa anotación de expertos es a la vez una limitación de la investigación y un recurso público que vale la pena construir; por ello, las figuras y tablas siguientes deben entenderse como evidencia humana valiosa pero preliminar, no como una garantía estadística que abarque todos los contextos clínicos.}

\begin{warningbox}{La identidad del evaluador determina qué preguntas puede responder}
Un paciente puede juzgar si una respuesta es comprensible y útil, pero no necesariamente si los detalles concuerdan con el consenso clínico; un médico puede juzgar el contenido médico, pero no necesariamente representa las necesidades de comunicación de todos los pacientes. La doble evaluación no consiste en que una sustituya a la otra, sino en colocar cada epistemic role en el lugar que le corresponde.
\end{warningbox}

\subsection{Resumen de la sección}

La evaluación de respuestas largas debe dividirse en varias dimensiones de riesgo y de utilidad, y debe distinguir las responsabilidades de juicio de los médicos y de los usuarios comunes. Este marco proporciona el sistema de coordenadas para los resultados de Med-PaLM que siguen, y muestra también que el cuello de botella de la investigación sobre modelos clínicos no son solo los datos de entrenamiento, sino también los datos de evaluación humana de alta calidad y verificables.
